\section{Results}
\label{sec:results}

Figure~\ref{fig:curves} shows the training loss of the three runs and
Table~\ref{tab:results} summarises them. The update rules of
Equations~\ref{eq:sgd}--\ref{eq:adam} differ most in how fast they make progress:
\nbval{fastest} crosses a training loss of $0.1$ after \nbval{fastest-epochs} epochs, and an
optimiser with ``--'' in the last column of the table never got there in some run.

\begin{figure}[ht]
  \centering
  \nbfig{fig-curves}
  \caption{Training loss per epoch on a logarithmic scale: mean over \nbval{n-seeds} seeds,
    with a band of one standard deviation. The dotted line marks $0.1$.}
  \label{fig:curves}
\end{figure}

\begin{table}[ht]
  \centering
  \nbtable{results}
  \caption{Final training loss, test accuracy, and the first epoch at which the training
    loss fell below $0.1$, as mean $\pm$ standard deviation over \nbval{n-seeds} seeds
    (``--'' if some run never got there).}
  \label{tab:results}
\end{table}

Test accuracy follows from that. The mean accuracies of the three optimisers span
\nbval{acc-spread} points, with \nbval{best} ahead at \nbval{best-acc}. Where the standard
deviations overlap the ranking should not be read as significant, and an optimiser that has
not finished converging is being compared on speed as much as on the quality of its solution.
Figure~\ref{fig:boundary} shows the decision boundary one of those runs learned.
% TODO: sweep the learning rate for each optimiser instead of fixing it by hand.
% TODO: add a second dataset; two moons alone says little about harder problems.

\begin{figure}[ht]
  \centering
  \nbfig{fig-boundary}
  \caption{Decision boundary of one \nbval{best} run, with the test points.}
  \label{fig:boundary}
\end{figure}
